\section{Configuración Experimental}

\subsection{Escenario}

Se desplegaron $N = 8$ robots en un arreglo en cuadrícula de aproximadamente $8\times9$~m. Cada robot recibe como objetivo el punto opuesto en la cuadrícula, forzando que todas las trayectorias converjan hacia la región central y generando situaciones de proximidad que activan el PVA. Las posiciones iniciales de los robots se muestran en la Fig.~\ref{fig:posiciones_iniciales} y las posiciones al final de la simulación en la Fig.~\ref{fig:posiciones_finales}.

\begin{figure}[H]
  \centering
  \includegraphics[width=0.9\columnwidth]{posiciones_iniciales}
  \caption{Posiciones iniciales de los 8 robots en la simulación.}
  \label{fig:posiciones_iniciales}
\end{figure}

\begin{figure}[H]
  \centering
  \includegraphics[width=0.9\columnwidth]{posiciones_finales}
  \caption{Posiciones finales de los robots al término de la simulación.}
  \label{fig:posiciones_finales}
\end{figure}

\subsection{Parámetros}

Los parámetros del sistema se resumen en la Tabla~\ref{tab:parametros}.

\begin{table}[H]
  \centering
  \caption{Parámetros de la simulación PVA}
  \label{tab:parametros}
  \begin{tabular}{lll}
    \toprule
    Parámetro & Símbolo & Valor \\
    \midrule
    Número de robots         & $N$              & 8 \\
    Distancia de seguridad   & $d_{\text{safe}}$  & 0.85~m \\
    Radio de influencia      & $d_{\text{inf}}$   & 3.0~m \\
    Parámetro de suavizado   & $\xi$            & 1.0 \\
    Radio del robot          & $r_p$            & 0.25~m \\
    Velocidad lineal máx.    & $v_{\max}$       & 1.0~m/s \\
    Velocidad angular máx.   & $\omega_{\max}$  & 1.0~rad/s \\
    Ganancia de control      & $k_1$            & 0.5 \\
    Ganancia de control      & $k_2$            & 1.0 \\
    Frecuencia de control    & $f_c$            & 20~Hz \\
    \bottomrule
  \end{tabular}
\end{table}
